% GENERATED FILE. Edit canonical lessons, then run npm run generate:books.
% canonical-source-hash: fnv1a64:20866f0106593f0b
% canonical-lessons: RU-C43-meat, RU-C43-fish, RU-C43-egg, RU-C43-butter, RU-C43-salt

\chapter{Meat, Fish, Egg, Butter, Salt}
\label{ch:ru-meat-fish-egg-butter-salt}
\hlchaptermodality{43}

\begin{chapteropening}
\textbf{By the end of this chapter:} \emph{I can name meat, fish, an egg, butter and salt.}

Complete the last lesson of chapter 43: I can name meat, fish, an egg, butter and salt.
\end{chapteropening}

\section[\texorpdfstring{\ru{мясо} (myáso)}{мясо (myáso)}]{\ru{мясо} (myáso) — meat}

\label{lesson:RU-C43-meat}



\begin{quote}
Before the new one: say the Russian for a question, then the Russian for an answer.
\end{quote}

\subsection*{You'll want to know: \ru{мясо}}
\textbf{\ru{мясо}} (\emph{myáso}) — \textquotedblleft{}meat\textquotedblright{}.

\begin{cousinweb}
\textbf{\ru{мясо}} $\leftarrow$ Proto-Slavic \emph{*męso} — inherited; the old nasal \emph{ę} became Russian \emph{\ru{я}}. Vasmer lists Old Prussian \emph{mensā}, Sanskrit \emph{māṃsá-} and Gothic \emph{mimz}, all \textquotedblleft{}meat\textquotedblright{}.
\end{cousinweb}

\begin{grammarlens}[title={What you've built}]
One more word for this chapter.
\end{grammarlens}

\subsection*{Your turn}
\begin{itemize}
\raggedright
  \item \emph{Say it:} \emph{myáso}
  \item \emph{Say it:} \emph{myáso}, once more
  \item \emph{Say it:} say \emph{myáso}
  \item \emph{Recall:} say the Russian for a question, then the Russian for an answer, then say \emph{myáso} again
\end{itemize}

\begin{tcolorbox}[breakable,colback=teal!4,colframe=teal!35!black,title={Before you move on}]
What does \ru{мясо} mean? (\textquotedblleft{}Meat\textquotedblright{}.) Say it once more.
\end{tcolorbox}
\section[\texorpdfstring{\ru{рыба} (rýba)}{рыба (rýba)}]{\ru{рыба} (rýba) — fish}

\label{lesson:RU-C43-fish}



\begin{quote}
Before the new one: say the Russian for an answer, then the Russian for meat.
\end{quote}

\subsection*{You'll want to know: \ru{рыба}}
\textbf{\ru{рыба}} (\emph{rýba}) — \textquotedblleft{}fish\textquotedblright{}.

\begin{cousinweb}
\textbf{\ru{рыба}} $\leftarrow$ Proto-Slavic \emph{*ryba} — inherited, but its deeper origin is unclear. Vasmer thought it was probably a taboo stand-in for an older word, the cousin of Lithuanian \emph{žuvìs}, \textquotedblleft{}fish\textquotedblright{}, which fishermen avoided because it sounded like \emph{\ru{звать}}, \textquotedblleft{}to call\textquotedblright{}. A comparison with Old High German \emph{rûppa}, \textquotedblleft{}caterpillar\textquotedblright{}, is also reported.
\end{cousinweb}

\begin{grammarlens}[title={What you've built}]
One more word for this chapter.
\end{grammarlens}

\subsection*{Your turn}
\begin{itemize}
\raggedright
  \item \emph{Say it:} \emph{rýba}
  \item \emph{Say it:} \emph{rýba}, once more
  \item \emph{Say it:} say \emph{rýba}
  \item \emph{Recall:} say the Russian for an answer, then the Russian for meat, then say \emph{rýba} again
\end{itemize}

\begin{tcolorbox}[breakable,colback=teal!4,colframe=teal!35!black,title={Before you move on}]
What does \ru{рыба} mean? (\textquotedblleft{}Fish\textquotedblright{}.) Say it once more.
\end{tcolorbox}
\section[\texorpdfstring{\ru{яйцо} (yaytsó)}{яйцо (yaytsó)}]{\ru{яйцо} (yaytsó) — an egg}

\label{lesson:RU-C43-egg}



\begin{quote}
Before the new one: say the Russian for meat, then the Russian for fish.
\end{quote}

\subsection*{You'll want to know: \ru{яйцо}}
\textbf{\ru{яйцо}} (\emph{yaytsó}) — \textquotedblleft{}an egg\textquotedblright{}.

\begin{cousinweb}
\textbf{\ru{яйцо}} $\leftarrow$ Proto-Slavic \emph{*aje}, \textquotedblleft{}egg\textquotedblright{}, with a diminutive ending — inherited. Vasmer lists Old High German \emph{ei}, Old Norse \emph{egg}, Latin \emph{ōvum} and Greek \emph{ōión}: English \emph{egg} is a cousin.
\end{cousinweb}

\begin{grammarlens}[title={What you've built}]
One more word for this chapter.
\end{grammarlens}

\subsection*{Your turn}
\begin{itemize}
\raggedright
  \item \emph{Say it:} \emph{yaytsó}
  \item \emph{Say it:} \emph{yaytsó}, once more
  \item \emph{Say it:} say \emph{yaytsó}
  \item \emph{Recall:} say the Russian for meat, then the Russian for fish, then say \emph{yaytsó} again
\end{itemize}

\begin{tcolorbox}[breakable,colback=teal!4,colframe=teal!35!black,title={Before you move on}]
What does \ru{яйцо} mean? (\textquotedblleft{}An egg\textquotedblright{}.) Say it once more.
\end{tcolorbox}
\section[\texorpdfstring{\ru{масло} (máslo)}{масло (máslo)}]{\ru{масло} (máslo) — butter; oil}

\label{lesson:RU-C43-butter}



\begin{quote}
Before the new one: say the Russian for fish, then the Russian for an egg.
\end{quote}

\subsection*{You'll want to know: \ru{масло}}
\textbf{\ru{масло}} (\emph{máslo}) — \textquotedblleft{}butter; oil\textquotedblright{}.

\begin{cousinweb}
\textbf{\ru{масло}} $\leftarrow$ Proto-Slavic \emph{*maslo}, from \emph{*maz-slo}, \textquotedblleft{}what you smear\textquotedblright{}, on the root of \emph{\ru{мазать}}, \textquotedblleft{}to smear\textquotedblright{} (Vasmer). So butter and oil share one word: both are spread.
\end{cousinweb}

\begin{grammarlens}[title={What you've built}]
One more word for this chapter.
\end{grammarlens}

\subsection*{Your turn}
\begin{itemize}
\raggedright
  \item \emph{Say it:} \emph{máslo}
  \item \emph{Say it:} \emph{máslo}, once more
  \item \emph{Say it:} say \emph{máslo}
  \item \emph{Recall:} say the Russian for fish, then the Russian for an egg, then say \emph{máslo} again
\end{itemize}

\begin{tcolorbox}[breakable,colback=teal!4,colframe=teal!35!black,title={Before you move on}]
What does \ru{масло} mean? (\textquotedblleft{}Butter; oil\textquotedblright{}.) Say it once more.
\end{tcolorbox}
\section[\texorpdfstring{\ru{соль} (sol)}{соль (sol)}]{\ru{соль} (sol) — salt}

\label{lesson:RU-C43-salt}



\begin{quote}
Before the new one: say the Russian for an egg, then the Russian for butter; oil.
\end{quote}

\subsection*{You'll want to know: \ru{соль}}
\textbf{\ru{соль}} (\emph{sol}) — \textquotedblleft{}salt\textquotedblright{}.

\begin{cousinweb}
\textbf{\ru{соль}} $\leftarrow$ Proto-Slavic \emph{*solĭ} — inherited. Vasmer lists Latvian \emph{sāls}, Latin \emph{sāl}, Greek \emph{háls} and Gothic \emph{salt}, all \textquotedblleft{}salt\textquotedblright{}: English \emph{salt} is a cousin.
\end{cousinweb}

\begin{grammarlens}[title={What you've built}]
One more word for this chapter.
\end{grammarlens}

\subsection*{Your turn}
\begin{itemize}
\raggedright
  \item \emph{Say it:} \emph{sol}
  \item \emph{Say it:} \emph{sol}, once more
  \item \emph{Say it:} say \emph{sol}
  \item \emph{Recall:} say the Russian for an egg, then the Russian for butter; oil, then say \emph{sol} again
\end{itemize}

\begin{tcolorbox}[breakable,colback=teal!4,colframe=teal!35!black,title={Before you move on}]
What does \ru{соль} mean? (\textquotedblleft{}Salt\textquotedblright{}.) Say it once more.
\end{tcolorbox}
